% Part: first-order-logic
% Chapter: natural-deduction
% Section: identity

\documentclass[../../../include/open-logic-section]{subfiles}

\begin{document}

\olfileid{fol}{ntd}{ide}

\olsection{\usetoken{P}{derivation} با \usetoken{S}{identity}}

!!^{derivation}s با !!{identity} به قواعد استنتاج افزوده‌ای نیاز دارند.

\begin{defish}
\AxiomC{}
\RightLabel{\Intro{\eq}}
\UnaryInfC{$\eq[t][t]$}
\DisplayProof
\hfill
\begin{tabular}{r}
\AxiomC{$\eq[t_1][t_2]$}
\AxiomC{$!A(t_1)$}
\RightLabel{\Elim{\eq}}
\BinaryInfC{$!A(t_2)$}
\DisplayProof
\\[3ex]
\AxiomC{$\eq[t_1][t_2]$}
\AxiomC{$!A(t_2)$}
\RightLabel{\Elim{\eq}}
\BinaryInfC{$!A(t_1)$}
\DisplayProof
\end{tabular}
\end{defish}

در قواعد بالا، $t$، $t_1$ و $t_2$ ترم‌های بسته‌اند. قاعدهٔ
\Intro{\eq} به ما اجازه می‌دهد هر گزارهٔ همانی به صورت
$\eq[t][t]$ را بی‌واسطه و بدون هیچ فرضی !!{derive} کنیم.

\begin{ex}
اگر $s$ و $t$ ترم‌های بسته باشند، آنگاه
$!A(s), \eq[s][t] \Proves !A(t)$:
\begin{prooftree}
\AxiomC{$\eq[s][t]$}
\AxiomC{$!A(s)$}
\RightLabel{$\Elim{\eq}$}
\BinaryInfC{$!A(t)$}
\end{prooftree}
شاید این اصل را با نام «اصل جانشینی همانی‌ها» یا قانون لایبنیتس
بشناسید.
\end{ex}

\begin{prob}
ثابت کنید $=$ هم متقارن و هم متعدی است؛ یعنی !!{derivation}sی از
$\lforall[x][\lforall[y][(\eq[x][y] \lif
    \eq[y][x])]]$ و $\lforall[x][\lforall[y][\lforall[z]((\eq[x][y]
    \land \eq[y][z]) \lif \eq[x][z])]]$ به دست دهید
\end{prob}

\begin{ex}
!!{derive} کردنِ !!{sentence} زیر را نشان می‌دهیم:
\begin{align*}
& \lforall[x][\lforall[y][((!A(x) \land !A(y)) \lif \eq[x][y])]]
\intertext{از !!{sentence}}
& \lexists[x][\lforall[y][(!A(y) \lif \eq[y][x])]]
\end{align*}
!!{derivation} را از پایان به آغاز می‌سازیم:
\begin{prooftree}
\AxiomC{$\lexists[x][\lforall[y][(!A(y) \lif \eq[y][x])]]
\quad \Discharge{!A(a) \land !A(b)}{1}$}
\DeduceC{$\eq[a][b]$}
\DischargeRule{\Intro{\lif}}{1}
\UnaryInfC{$((!A(a) \land !A(b)) \lif \eq[a][b])$}
\RightLabel{\Intro{\lforall}}
\UnaryInfC{$\lforall[y][((!A(a) \land !A(y)) \lif \eq[a][y])]$}
\RightLabel{\Intro{\lforall}}
\UnaryInfC{$\lforall[x][\lforall[y][((!A(x) \land !A(y)) \lif \eq[x][y])]]$}
\end{prooftree}
اکنون باید از فرض اصلی استفاده کنیم: چون این فرض یک !!{formula}
وجودی است، از \Elim{\lexists} استفاده می‌کنیم تا نتیجهٔ میانی
$\eq[a][b]$ را !!{derive} کنیم.
\begin{prooftree}
\AxiomC{$\lexists[x][\lforall[y][(!A(y) \lif \eq[y][x])]]$}
\AxiomC{$\Discharge{\lforall[y][(!A(y) \lif \eq[y][c]])}{2}$}
\noLine
\UnaryInfC{$\Discharge{!A(a) \land !A(b)}{1}$}
\DeduceC{$\eq[a][b]$}
\DischargeRule{\Elim{\lexists}}{2}
\BinaryInfC{$\eq[a][b]$}
\DischargeRule{\Intro{\lif}}{1}
\UnaryInfC{$((!A(a) \land !A(b)) \lif \eq[a][b])$}
\RightLabel{\Intro{\lforall}}
\UnaryInfC{$\lforall[y][((!A(a) \land !A(y)) \lif \eq[a][y])]$}
\RightLabel{\Intro{\lforall}}
\UnaryInfC{$\lforall[x][\lforall[y][((!A(x) \land !A(y)) \lif \eq[x][y])]]$}
\end{prooftree}
!!{derivation} فرعیِ بالا سمت راست با استفاده از فرض‌هایش برای نشان‌دادن
$\eq[a][c]$ و $\eq[b][c]$ کامل می‌شود. این امر به دو !!{derivation}s
جداگانه نیاز دارد. !!{derivation} برای $\eq[a][c]$ چنین است:
\begin{prooftree}
\AxiomC{$\Discharge{\lforall[y][(!A(y) \lif \eq[y][c]])}{2}$}
\RightLabel{\Elim{\lforall}}
\UnaryInfC{$!A(a) \lif \eq[a][c]$}
\AxiomC{$\Discharge{!A(a) \land !A(b)}{1}$}
\RightLabel{\Elim{\land}}
\UnaryInfC{$!A(a)$}
\RightLabel{\Elim{\lif}}
\BinaryInfC{$\eq[a][c]$}
\end{prooftree}
از $\eq[a][c]$ و $\eq[b][c]$، !!{derive} کردنِ $\eq[a][b]$ با
\Elim{\eq} امکان‌پذیر است.
\end{ex}

\begin{prob}
!!{derivation}sی از !!{formula}s زیر به دست دهید:
\begin{enumerate}
\item $\lforall[x][\lforall[y][((\eq[x][y] \land !A(x)) \lif !A(y))]]$
\item $\lexists[x][!A(x)] \land \lforall[y][\lforall[z][((!A(y) \land
    !A(z)) \lif \eq[y][z])]] \lif \lexists[x][(!A(x) \land
  \lforall[y][(!A(y) \lif \eq[y][x])])]$
\end{enumerate}
\end{prob}

\end{document}
